% PAPER_TEMPLATE.tex
% Comprehensive LaTeX paper template for CA29 (SAC implementation)
% Usage: pdflatex PAPER_TEMPLATE.tex && bibtex PAPER_TEMPLATE && pdflatex PAPER_TEMPLATE.tex && pdflatex PAPER_TEMPLATE.tex

\documentclass[11pt,a4paper]{article}
\usepackage[utf8]{inputenc}
\usepackage{amsmath,amssymb,amsthm}
\usepackage{graphicx}
\usepackage{booktabs}
\usepackage{hyperref}
\usepackage{cleveref}
\usepackage{siunitx}
\usepackage{caption}
\usepackage{subcaption}
\usepackage{geometry}
\geometry{margin=1in}
\usepackage{natbib}
\usepackage{algorithm}
\usepackage{algpseudocode}
\usepackage{float}
\usepackage{listings}
\usepackage{tcolorbox}
\setcitestyle{numbers,comma,square}

\title{Soft Actor-Critic (SAC) for Continuous Control: Implementation, Evaluation, and Reproducibility}
\author{Author Name \\\small Affiliation}
\date{\today}

\begin{document}
\maketitle

\begin{abstract}
We provide a comprehensive implementation of the Soft Actor-Critic (SAC) algorithm and an experimental pipeline designed for reproducibility and rigorous evaluation. This document contains algorithmic derivations, detailed implementation notes to ensure numerical stability, a complete experimental protocol (including ablations), and templates for reporting and reproducibility. The repository includes unit tests, a CLI, notebook-based visualizations, and LaTeX templates for producing a formal paper.
\end{abstract}

\section{Introduction}
Reinforcement learning for continuous control requires balancing stability and exploration. Soft Actor-Critic (SAC) \citep{haarnoja2018soft} achieves this by adding an entropy bonus to the return, promoting exploratory stochastic policies while maintaining sample efficiency via off-policy learning. This work documents our implementation choices and provides tools and templates for reproducible research.

\subsection{Contributions}
\begin{itemize}
  \item A robust, import-safe PyTorch implementation of SAC with unit tests and a simple CLI.
  \item An experiment runner (`Experiment`) that handles seeding, evaluation, logging, and metric collection into `metrics.csv`.
  \item Notebooks and LaTeX templates for quick analysis and a paper-style writeup, respectively.
\end{itemize}

\section{Related Work}
SAC builds on deterministic and stochastic actor-critic families (DDPG \citep{lillicrap2015continuous}, TD3 \citep{fujimoto2018addressing}) and connects to maximum entropy RL literature. Our implementation borrows best practices from recent reproducible RL work, focusing on numerical stability and consistent logging for experiments.

\section{Method}
\subsection{Maximum Entropy Objective and Notation}
Given state space $\mathcal{S}$ and action space $\mathcal{A}$, SAC optimizes
\begin{equation}
J(\pi) = \E_{\tau\sim\pi} \left[ \sum_{t=0}^{\infty} r(s_t,a_t) + \alpha \mathcal{H}(\pi(\cdot|s_t)) \right],
\end{equation}
where $\mathcal{H}(\pi(\cdot|s)) = -\E_{a\sim\pi}[\log \pi(a|s)]$.

\subsection{Policy Representation}
We use a squashed Gaussian policy: sample $x\sim\mathcal{N}(\mu_\theta(s),\sigma_\theta(s))$, set $a=\tanh(x)$, and compute corrected log-probabilities:
\begin{equation}
\log \pi(a|s) = \log \mathcal{N}(x;\mu,\sigma) - \sum_j \log \left(1 - \tanh(x_j)^2 + \epsilon\right).
\end{equation}
This correction is critical for correct entropy estimation and stable gradients.

\subsection{Critic (Q) Learning}
We train two critics $Q_{\phi_1}$ and $Q_{\phi_2}$, and the target is:
\begin{equation}
y = r + \gamma \left(\min_{i=1,2} Q_{\phi'_i}(s', a') - \alpha \log \pi_\theta(a'|s') \right).
\end{equation}
Minimizing $(Q_\phi(s,a) - y)^2$ yields stable critic updates; twin critics reduce overestimation.

\subsection{Automatic Temperature Tuning}
We optionally adapt $\alpha$ to maintain target entropy $\mathcal{H}_0$ by minimizing
\begin{equation}
\mathcal{L}(\alpha) = -\E_{a\sim\pi_\theta}[\alpha(\log\pi_\theta(a|s) + \mathcal{H}_0)].
\end{equation}
Using the log parameterization $\log \alpha$ helps constrain positivity and stabilizes optimization.

\section{Algorithm}
\begin{algorithm}[H]
\caption{Soft Actor-Critic (detailed)}
\begin{algorithmic}[1]
\State Initialize $\pi_\theta$, $Q_{\phi_1}$, $Q_{\phi_2}$, target networks $\phi'_1,\phi'_2$, replay buffer $\mathcal{D}$
\State Set target entropy $\mathcal{H}_0$ (default $-|\mathcal{A}|$) and initial log alpha
\For{each environment step}
  \State Collect $(s,a,r,s',d)$ using current policy
  \State Add to $\mathcal{D}$
  \If{time to update}
    \For{gradient steps}
      \State Sample minibatch from $\mathcal{D}$
      \State Compute targets $y$ and critic losses; update critics
      \State Compute actor loss using sampled actions; update policy
      \If{automatic entropy tuning} update $\alpha$ via gradient on dual objective
      \State Soft-update targets: $\phi' \leftarrow \tau \phi + (1-\tau)\phi'$
    \EndFor
  \EndIf
\EndFor
\end{algorithmic}
\end{algorithm}

\section{Implementation Notes and Best Practices}
\subsection{Numerical Stability}
- Add small constant $\epsilon$ when computing log or dividing.
- Clamp `log_std` outputs (e.g., [-20,2]) to keep std in a reasonable range.
- Use double precision where appropriate for numerical checks; training in fp32 is standard.

\subsection{Seeding and Determinism}
Use `set_seed` and `set_env_seed` utilities. Note: deterministic execution can still vary across hardware and PyTorch versions.

\section{Experiments}
\subsection{Environments}
We recommend HalfCheetah-v4, Walker2d-v4, Hopper-v4; for more computationally expensive tests use Ant-v4.

\subsection{Evaluation Protocol}
Run each config with multiple seeds; evaluate every `eval_freq` steps for `num_eval_episodes` episodes. Store per-evaluation mean and std in `metrics.csv` for plotting.

\subsection{Ablation Design}
Fix a single variable (e.g., alpha, buffer size, network width) and run a sweep. Keep seeds consistent across the sweep for fair comparisons.

\section{Results Reporting Format}
Include the following in the `results/<tag>/` folder:
\begin{itemize}
  \item `metrics.csv` (columns: step, avg_reward, std_reward)
  \item `training.log` (logger output)
  \item `sac_final.pth` (checkpoint)
  \item `plots/` (PNG/PDF for figures referenced in the paper)
  \item `metadata.txt` (git sha, config used, seeds, environment versions)
\end{itemize}

Include a LaTeX table with final results per environment (example in the file) and embed plots created from `metrics.csv`.

\section{Discussion and Future Work}
Discuss algorithmic strengths and failure modes observed, propose extensions (recurrent policies, distributional critics, improved exploration priors), and suggest additional stress tests.

\section{Reproducibility Appendix}
Provide exact commands used to run important experiments. Example:
\begin{lstlisting}[language=bash]
python -m src.cli --config configs/default.yaml --log-dir results/halfcheetah_default

git rev-parse HEAD > results/halfcheetah_default/commit.txt
\end{lstlisting}

\section*{Acknowledgements}
We acknowledge contributors and reviewers and any funding sources.

\bibliographystyle{abbrvnat}
\bibliography{references}

\appendix
\section{Appendix A: Hyperparameter Table}
\begin{table}[H]
\centering
\caption{Primary hyperparameters used in experiments and suggested ranges}
\begin{tabular}{lcc}
\toprule
Parameter & Default & Suggested range \\
\midrule
$\gamma$ & 0.99 & 0.95 -- 0.999 \\
actor lr & 3e-4 & 1e-4 -- 1e-3 \\
critic lr & 3e-4 & 1e-4 -- 1e-3 \\
batch size & 256 & 64 -- 512 \\
buffer size & 1e6 & 1e5 -- 1e7 \\
$\tau$ & 0.005 & 0.001 -- 0.01 \\
initial $\alpha$ & 0.2 & 1e-4 -- 1.0 \\
\bottomrule
\end{tabular}
\end{table}

\section{Appendix B: Compute and Runtime}
Document sample runtime and GPU usage for representative runs here (e.g., 100k steps takes X hours on Y GPU with Z memory).

\section{Appendix C: Additional Plots and Seed-level Results}
Include extra plots (per-seed curves, variance breakdowns, entropy traces) here or in the supplemental materials.

\end{document}
